problem:  
Let \( (M^n,g) \) be a compact, simply connected Riemannian manifold with **quasipositive curvature** (\(\sec \ge 0\) everywhere and \(\sec > 0\) somewhere).  Suppose \(M\) admits an **effective cohomogeneity‑one** isometric action by a compact connected Lie group \(G\), determined by a group diagram \((G,H_-,H_+,K)\).  

In known classifications of positively curved cohomogeneity‑one manifolds, the existence of an invariant metric with \(\sec>0\) is an extremely restrictive property.  Certain diagrams (for example, some involving unequal isotropy slice dimensions or curvature‑sign changes along normal geodesics) are widely believed or proven **not** to support any \(G\)-invariant positively curved metric, even though they permit nonnegatively curved ones.  

**Question:**  
Does there exist a compact, simply connected, cohomogeneity‑one manifold \(M\) with a smooth \(G\)-invariant quasipositively curved metric such that **no \(G\)-invariant positively curved metric exists on \(M\)**?  
Equivalently, can quasipositive curvature appear on cohomogeneity‑one manifolds that are intrinsically excluded from positive curvature within the same symmetry class?

Why is it a "good" problem:  
- **Profound insight:**  
  This formulation asks whether quasipositivity introduces fundamentally new geometric behavior even under high symmetry.  It tests whether the step from \(\sec>0\) to \(\sec\ge0\) with some positivity has real structural consequences or whether curvature rigidity persists unchanged.  Understanding this transition would illuminate deep mechanisms controlling how local zero‑curvature regions interact with global orbit geometry.  

- **Bridge between fields:**  
  Resolving this problem would require integrating techniques from **Riemannian geometry** (analysis of curvature for invariant metrics), **representation theory** (study of isotropy actions), and **topology of transformation groups** (cohomogeneity‑one gluing constructions).  It could also connect to **metric collapse** and **Alexandrov geometry**, exploring how curvature degenerations manifest under symmetry.  

- **Extreme simplicity and depth:**  
  The statement—whether a quasipositively curved, \(G\)-invariant metric can exist where every \(G\)-invariant positively curved metric is excluded—is transparent but reaches into delicate geometric territory.  A single explicit counterexample or proof of impossibility would have strong conceptual impact.  

- **Potential impact:**  
  A positive discovery would produce the first instance of a truly *new* cohomogeneity‑one manifold type arising only in the quasipositive regime, expanding the catalog of known curvature geometries.  A negative result would unify both curvature categories under a single rigidity principle.  Either conclusion would significantly clarify the boundary between nonnegative, quasipositive, and positive curvature in the presence of symmetry.  

Therefore, this problem formulates a **precise, intrinsic, and mathematically sound** challenge lying exactly on the frontier between symmetry‑constrained curvature geometry and topological rigidity, making it a genuinely **good research problem**.